% =====================================================================
%  Distributed Multi-Hop Indoor Localization via Transformer-Based
%  Joint Distance, Angle, and Confidence Estimation from Wi-Fi CSI
%
%  IEEE-style paper (conference format, IEEEtran)
%  Generated from the codebase:
%    - train_transformer_dist_angle_conf.py
%    - localization_transformer.py
%  Compile with pdflatex/bibtex (or simply pdflatex with inline thebibliography).
% =====================================================================

\documentclass[conference]{IEEEtran}
\usepackage[utf8]{inputenc}
\usepackage{cite}
\usepackage{amsmath,amssymb,amsfonts}
\usepackage{graphicx}
\usepackage{textcomp}
\usepackage{xcolor}
\usepackage{booktabs}
\usepackage{multirow}
\usepackage{balance}

\def\BibTeX{{\rm B\kern-.05em{\sc i\kern-.025em b}\kern-.08em
    T\kern-.1667em\lower.7ex\hbox{E}\kern-.125emX}}

\begin{document}

\title{Distributed Multi-Hop Indoor Localization via Transformer-Based Joint Distance, Angle, and Confidence Estimation from Wi-Fi CSI}

\author{\IEEEauthorblockN{Author Name$^{1}$ and Author Name$^{2}$}
\IEEEauthorblockA{\textit{$^{1}$Department, University}, City, Country\\
$^{1}$email@example.com}
\IEEEauthorblockA{\textit{$^{2}$Department, University}, City, Country\\
$^{2}$email@example.com}
}

% NOTE: Replace the placeholder author block with the actual authors
% and affiliations before submission.

\maketitle

\begin{abstract}
Indoor localization using Wi-Fi Channel State Information (CSI) is an attractive
alternative to satellite positioning, yet existing learning-based approaches
typically regress a single geometric quantity and assume line-of-sight (LOS)
propagation. We present an end-to-end framework that couples a
\textit{Transformer} jointly estimating \textbf{distance}, \textbf{angle}, and
\textbf{confidence} from OFDM CSI with a \textit{distributed multi-hop}
localization algorithm. The Transformer treats every subcarrier as a token,
appends a learnable class (CLS) token, and routes the encoded representation to
three task-specific multi-layer perceptron (MLP) heads. A confidence-weighted
multi-task loss scales the angle objective by the true propagation confidence,
so that non-line-of-sight (NLOS) samples whose angle is physically unlearnable
do not dominate the gradient. The estimated confidence then acts as a soft
LOS/NLOS gate: a base station (BS) iteratively localizes the neighbour with the
highest predicted confidence, and each newly localized node propagates the
process outward, forming a multi-hop localization tree. Finally, a
confidence-weighted graph optimization refines all positions jointly to
suppress accumulated drift. On a simulated 5.8\,GHz, 20\,MHz, 4-antenna indoor
dataset spanning 25 distances, 36 angles, and 10 wall-obstruction levels, the
proposed model achieves a distance MAE of 0.81\,m and an angle MAE of 8.2$^\circ$
in LOS conditions, while the confidence head cleanly separates LOS from NLOS
links (0.75 vs.\ 0.35 mean confidence), enabling robust distributed
localization over a 16-node, wall-partitioned floor plan.
\end{abstract}

\begin{IEEEkeywords}
channel state information, indoor localization, transformer, angle of arrival,
distance estimation, multi-hop localization, deep learning
\end{IEEEkeywords}

\section{Introduction}
\IEEEPARstart{I}{ndoor} localization is a cornerstone of context-aware
applications such as asset tracking, autonomous navigation, and the Internet of
Things (IoT). Global Navigation Satellite Systems (GNSS) fail indoors owing to
strong signal attenuation and multipath, motivating the use of ubiquitous
Wi-Fi infrastructure as a positioning source. Wi-Fi Channel State Information
(CSI) describes the frequency-selective channel between a transmitter and a
receiver, and---unlike received signal strength (RSS)---it retains fine-grained
phase and amplitude structure that can be exploited for both ranging and
angle-of-arrival (AoA) estimation~\cite{halperin2011tool}.

Two families of approaches dominate the literature. \emph{Model-based} methods
such as SpotFi~\cite{kotaru2015spotfi} and ArrayTrack~\cite{xiong2013arraytrack}
employ super-resolution spectral estimation (e.g., MUSIC) on the CSI matrix to
estimate AoA and time-of-flight (ToF). \emph{Learning-based} methods such as
DeepFi~\cite{wang2015deepfi}, PhaseFi~\cite{wang2016phasefi}, and
ConFi~\cite{chen2017confi} feed hand-crafted or raw CSI into neural networks to
regress position. However, most prior work suffers from three limitations.
First, distance and angle are usually estimated by \emph{separate} models or
pipeline stages, ignoring their physical coupling. Second, NLOS propagation---a
pervasive condition indoors due to walls---is either ignored or handled with ad
hoc heuristics, even though the AoA estimate becomes essentially meaningless
when the dominant path is obstructed. Third, localization is typically
\emph{centralized}: a single AP must observe every node directly, which does not
scale to multi-room deployments where many nodes are hidden behind walls from
the AP.

This paper addresses these limitations with two coupled contributions:

\begin{itemize}
\item \textbf{A single-tower Transformer} that jointly regresses
      \emph{distance}, \emph{angle}, and a \emph{confidence} score in $[0,1]$
      from per-subcarrier CSI tokens. A confidence-weighted multi-task loss
      scales the angle objective by the true confidence, forcing the model to
      focus its angle-learning capacity on LOS samples while learning to emit
      low confidence for NLOS samples.
\item \textbf{A distributed multi-hop localization algorithm} in which the BS
      greedily localizes the highest-confidence neighbour, and each newly
      localized node recursively localizes its own neighbours. The learned
      confidence acts as a soft LOS/NLOS gate that determines the order and
      path of the localization tree, followed by a confidence-weighted graph
      optimization that refines all positions jointly.
\end{itemize}

The remainder of the paper is organized as follows. Section~\ref{sec:model}
describes the CSI signal model and data generation. Section~\ref{sec:features}
presents the CSI feature extraction. Section~\ref{sec:transformer} details the
Transformer estimator and loss. Section~\ref{sec:localization} describes the
distributed localization algorithm and Section~\ref{sec:graph} the graph
refinement. Section~\ref{sec:results} reports experimental results, and
Section~\ref{sec:conclusion} concludes.

\section{System Model and CSI Generation}
\label{sec:model}

We consider an indoor OFDM system operating at a carrier frequency
$f_c=5.8$\,GHz with a bandwidth of $20$\,MHz, $N_{\mathrm{fft}}=64$ subcarriers
(all used), a cyclic-prefix length $N_{\mathrm{cp}}=16$, and a subcarrier
spacing $\Delta f = F_s/N_{\mathrm{fft}} = 312.5$\,kHz. The receiver is equipped
with a $2\times 2$ rectangular antenna array of $N_{\mathrm{rx}}=4$ elements
spaced by $d_{\mathrm{ant}}=\lambda/3 \approx 1.72$\,cm, where
$\lambda = c/f_c \approx 5.17$\,cm. The array positions are
$\{(0,0),(d_{\mathrm{ant}},0),(0,d_{\mathrm{ant}}),(d_{\mathrm{ant}},d_{\mathrm{ant}})\}$,
so that the $x$-axis and $y$-axis antenna pairs encode $\cos\theta$ and
$\sin\theta$ respectively, resolving AoA unambiguously over the full
$360^{\circ}$~\cite{xiong2013arraytrack}.

Each link is modelled as a single-path Rayleigh channel with free-space path
loss (FSPL), wall penetration loss, and additive white Gaussian noise (AWGN).
For a transmitter--receiver pair at distance $d$ and angle $\theta$ (measured
counter-clockwise from the positive $x$-axis), the frequency-domain channel at
receive antenna $r$ and subcarrier $k$ is
\begin{equation}
  H_{r,k} = h\, g \,
    e^{-j 2\pi k \Delta f \tau}
    e^{-j 2\pi\left(x_r \cos\theta + y_r \sin\theta\right)/\lambda},
\label{eq:channel}
\end{equation}
where $h \sim \mathcal{CN}(0,1)$ is a per-packet Rayleigh fading coefficient,
$\tau = d/c$ is the propagation delay, $(x_r,y_r)$ is the position of antenna
$r$, and $g$ is the path gain. The phase rotation over subcarriers encodes the
ToF (hence distance), while the inter-antenna phase encodes the AoA. The path
gain combines FSPL and wall loss,
\begin{equation}
  g = 10^{-\left(\mathrm{FSPL} + L_{\mathrm{wall}}\right)/20},
  \quad
  \mathrm{FSPL} = 20\log_{10}\!\left(\frac{4\pi d}{\lambda}\right),
\label{eq:pathgain}
\end{equation}
where $L_{\mathrm{wall}} = \sum_{w} t_w \cdot a$ is the total penetration loss
over all crossed walls, with wall thickness $t_w$ and attenuation coefficient
$a = 100$\,dB/m. The reference SNR (no-wall baseline) is $25$\,dB; a fixed
noise floor referenced to the no-wall FSPL means that wall loss directly reduces
the effective SNR.

QPSK pilot symbols $s_k$ are transmitted, and the CSI is estimated at the
receiver by zero-forcing:
\begin{equation}
  \hat{H}_{r,k} = \frac{y_{r,k}}{s_k} = H_{r,k} + n_{r,k},
\label{eq:csi}
\end{equation}
where $n_{r,k} \sim \mathcal{CN}(0,\sigma_n^2)$ is the estimation noise. Each
link yields a complex CSI tensor $\hat{\mathbf{H}} \in
\mathbb{C}^{N_{\mathrm{pkt}}\times N_{\mathrm{rx}}\times N_{\mathrm{fft}}}$,
with $N_{\mathrm{pkt}}=500$ packets per link.

\subsection{Training Dataset}
Ten obstruction levels are generated with wall losses
$L_{\mathrm{wall}} \in \{0,10,\dots,90\}$\,dB, corresponding to confidence
labels $c \in \{1.0,0.9,\dots,0.1\}$. For each level, a grid of $25$ distances
($1$--$49$\,m in $2$\,m steps) and $36$ angles ($0^{\circ}$--$350^{\circ}$ in
$10^{\circ}$ steps) is simulated with $500$ packets per combination, giving
$900$ combinations and $450{,}000$ samples per level. Samples with distance
$>30$\,m are discarded to reflect realistic Wi-Fi coverage, leaving
$270{,}000$ samples per level ($2.7$\,M total), subsampled to $500{,}000$
samples for training.

\section{CSI Feature Extraction}
\label{sec:features}

Raw CSI is transformed into a fixed-size real feature vector of dimension
$D = 2N_{\mathrm{fft}} + 2(N_{\mathrm{rx}}-1)N_{\mathrm{fft}} = 512$ per packet,
concatenating two complementary descriptors.

\paragraph{Distance feature.}
The complex CSI of the reference antenna (antenna 0) is flattened into its real
and imaginary parts:
\begin{equation}
  \mathbf{f}_{\mathrm{dist}} =
  \big[\Re\{\hat{H}_{0,:}\},\; \Im\{\hat{H}_{0,:}\}\big] \in \mathbb{R}^{2N_{\mathrm{fft}}},
\label{eq:distfeat}
\end{equation}
which preserves the subcarrier phase ramp induced by the ToF.

\paragraph{Angle feature.}
The phase of each antenna is first unwrapped along the subcarrier axis. The
inter-antenna phase differences relative to antenna 0 are then computed and
encoded as sine/cosine pairs to avoid the $2\pi$ wrap-around discontinuity:
\begin{equation}
  \Delta\phi_{r,k} = \mathrm{unwrap}\left[\angle\hat{H}_{r,k}\right]
                    - \mathrm{unwrap}\left[\angle\hat{H}_{0,k}\right],
  \qquad r=1,2,3,
\label{eq:phasediff}
\end{equation}
\begin{equation}
  \mathbf{f}_{\mathrm{ang}} = \big[\sin(\Delta\phi_{r,:}),\;
                              \cos(\Delta\phi_{r,:})\big]_{r=1}^{3}
                              \in \mathbb{R}^{6N_{\mathrm{fft}}},
\label{eq:angfeat}
\end{equation}
where the $x$-axis pair ($r=1$), $y$-axis pair ($r=2$), and diagonal pair
($r=3$) encode $\cos\theta$, $\sin\theta$, and $\cos\theta+\sin\theta$,
respectively. The final feature vector is
$\mathbf{f} = [\mathbf{f}_{\mathrm{dist}}; \mathbf{f}_{\mathrm{ang}}]$.

\section{Transformer-Based Joint Estimator}
\label{sec:transformer}

\subsection{Architecture}
The feature vector $\mathbf{f}$ is reshaped into $N_{\mathrm{fft}}=64$
per-subcarrier tokens, each of dimension
$d_{\mathrm{per}} = 2 + 2(N_{\mathrm{rx}}-1) = 8$ (antenna-0 real/imag plus
six phase-difference sine/cosine channels). As illustrated in
Fig.~\ref{fig:arch}, a linear projection maps each token to a model dimension
$d_{\mathrm{model}}=64$, a learnable positional embedding is added, and a
learnable class (CLS) token is prepended. A standard Transformer encoder with
$L=4$ layers, $H=4$ attention heads, feed-forward dimension $d_{\mathrm{ff}}=256$,
GELU activations, and dropout $0.1$ processes the sequence. The CLS output is
routed to three independent two-layer MLP heads (hidden sizes $128$--$64$):
a \emph{distance} head and an \emph{angle} head (both linear), and a
\emph{confidence} head with a sigmoid activation restricting the output to
$[0,1]$:
\begin{equation}
  [\hat{d},\; \hat{\theta},\; \hat{c}] = \mathrm{MLP}_{d}\oplus
    \mathrm{MLP}_{\theta}\oplus \mathrm{MLP}_{c}
    \left(\mathrm{CLS}\!\left[\mathrm{Encoder}(\mathbf{Z})\right]\right).
\label{eq:heads}
\end{equation}
The total model size is $254{,}595$ parameters ($0.97$\,MB), making it
deployable on resource-constrained edge devices.

\begin{figure}[t]
\centering
\fbox{\parbox{0.95\columnwidth}{\centering
  \vspace{4mm}
  \textbf{Per-subcarrier tokens}\\
  $\mathbf{f} \to \mathbb{R}^{64 \times 8}$\\
  \vspace{2mm}
  $\downarrow$ Linear projection + positional embedding\\
  $\downarrow$ prepend CLS token\\
  \vspace{2mm}
  \textbf{Transformer Encoder}\\
  ($L{=}4$ layers, $H{=}4$ heads, $d_{\mathrm{model}}{=}64$)\\
  \vspace{2mm}
  $\downarrow$ CLS output\\
  \vspace{2mm}
  \begin{tabular}{ccc}
   MLP$_{d}$ & MLP$_{\theta}$ & MLP$_{c}$\\
   $\hat{d}$ & $\hat{\theta}$ & $\hat{c}\in[0,1]$
  \end{tabular}\\
  \vspace{4mm}
}}
\caption{Transformer architecture for joint distance, angle, and confidence
estimation.}
\label{fig:arch}
\end{figure}

\subsection{Confidence-Weighted Multi-Task Loss}
The three heads are trained with a weighted mean-squared-error (MSE) loss. A key
design choice is to scale the angle term by the \emph{true} confidence label
$c$, so that NLOS samples---whose AoA is physically unlearnable under severe
obstruction---are down-weighted rather than corrupting the angle gradient:
\begin{equation}
  \mathcal{L} = w_{d}\,\mathbb{E}\!\left[(\hat{d}-d)^2\right]
    + w_{\theta}\,\mathbb{E}\!\left[c\,(\hat{\theta}-\theta)^2\right]
    + w_{c}\,\mathbb{E}\!\left[(\hat{c}-c)^2\right],
\label{eq:loss}
\end{equation}
with weights $w_{d}=1.0$, $w_{\theta}=5.0$, and $w_{c}=2.0$ chosen so that the
three terms contribute comparably. Distance and angle are standardized before
the loss, while confidence is regressed directly in $[0,1]$.

\section{Distributed Multi-Hop Localization}
\label{sec:localization}

The localization stage operates over a floor plan of $N$ nodes---one BS and
$N-1$ user equipment (UE) nodes---partitioned by walls. Each node is assigned an
IP address and a binary flag (0 = unlocalized, 1 = localized). The BS, whose
position is known a priori, is initialized with flag 1. The procedure, which
generalizes the anchor-based scheme of prior work, is as follows.

\begin{enumerate}
\item \textbf{Per-link estimation.} Each already-localized node collects
      $N_{\mathrm{pkt}}$ packets from each unlocalized neighbour and runs the
      Transformer to obtain per-packet estimates
      $(\hat{d},\hat{\theta},\hat{c})$, averaged over packets to yield
      $(\bar{d},\bar{\theta},\bar{c})$.
\item \textbf{Confidence-greedy selection.} Among all
      (localized $\to$ unlocalized) links, the link with the highest mean
      confidence $\bar{c}$ is selected, defining the next anchor and target.
\item \textbf{Position propagation.} The target's absolute position is computed
      from the anchor's known position and the estimated relative geometry:
      \begin{equation}
        \mathbf{p}_{\mathrm{target}} = \mathbf{p}_{\mathrm{anchor}}
          + \bar{d}\,[\cos\bar{\theta},\; \sin\bar{\theta}]^{\top},
      \label{eq:position}
      \end{equation}
      and its flag is set to 1.
\item \textbf{Iteration.} Steps 1--3 repeat until every node has flag 1,
      producing a spanning localization tree rooted at the BS.
\end{enumerate}

The learned confidence plays a crucial role: it prioritizes reliable LOS links
early in the process, preventing error propagation along obstructed links. If
the selected path becomes disconnected (no reachable unlocalized node), the
algorithm terminates gracefully with a partial solution.

\section{Graph-Based Refinement}
\label{sec:graph}

The sequential procedure localizes each node from a \emph{single} link, so
errors accumulate along the tree. To mitigate this drift, a global refinement
step minimizes the discrepancy between all Transformer-estimated link
measurements and those implied by the estimated positions. Let
$\mathcal{E}$ be the set of links retained after a confidence threshold
$\bar{c} \ge 0.8$ (the LOS-only links). The optimization solves
\begin{equation}
  \min_{\{\mathbf{p}_i\}}
  \sum_{(i,j)\in\mathcal{E}} w_{ij}
    \left( \|\mathbf{p}_i-\mathbf{p}_j\| - \bar{d}_{ij} \right)^2
  + \lambda_{\theta}\, w_{ij}
    \left( \|\mathbf{p}_i-\mathbf{p}_j\|\,\Delta\Theta_{ij} \right)^2,
\label{eq:graphcost}
\end{equation}
where $w_{ij} = \bar{c}_{ij}^{\,2}$ weights each link by its squared confidence,
$\Delta\Theta_{ij}$ is the wrapped angular error between the estimated angle and
the angle implied by the current positions, and $\lambda_{\theta}$ trades off
distance and angle fidelity. The lateral term
$\|\mathbf{p}_i-\mathbf{p}_j\|\Delta\Theta_{ij}$ scales the angular error so
that a $1^{\circ}$ error costs equally at any distance. The BS position is held
fixed as an anchor, and the problem is solved with the L-BFGS-B algorithm.

\section{Experimental Evaluation}
\label{sec:results}

\subsection{Setup}
The model is trained on $500{,}000$ samples split into train/validation/test
sets of $350{,}000$/$50{,}000$/$100{,}000$. Optimization uses AdamW
($\mathrm{lr}=3\times10^{-4}$) with a cosine-annealing warm-restart schedule
($T_0=30$, $T_{\mathrm{mult}}=2$, $\eta_{\min}=10^{-6}$), batch size $516$, and
early stopping. The best checkpoint is selected by validation loss (epoch $88$,
$\mathrm{val\ loss}=1.128$). Evaluation is performed on a separate held-out
test set of $1{,}080{,}000$ samples drawn from all ten obstruction levels, with
a held-out distance/angle grid ($2$--$50$\,m, $5^{\circ}$--$355^{\circ}$).

\subsection{Joint Estimation Accuracy}
Table~\ref{tab:overall} reports the overall test performance. Distance
estimation is robust across propagation conditions (MAE $0.81$\,m), whereas
angle estimation is excellent in LOS but degrades sharply in NLOS. The
confidence head cleanly separates the two regimes, assigning mean confidence
$0.75$ to LOS and $0.35$ to NLOS samples.

\begin{table}[t]
\centering
\caption{Overall test results (1,080,000 samples).}
\label{tab:overall}
\begin{tabular}{lccc}
\toprule
\multirow{2}{*}{Quantity} & \multicolumn{2}{c}{Overall} & \multirow{2}{*}{LOS vs.\ NLOS}\\
 & MAE & RMSE & \\
\midrule
Distance (m)   & 0.807 & 1.103 & 0.727 / 0.887 \\
Angle ($^{\circ}$) & 56.97 & 84.62 & 23.57 / 90.37 \\
Confidence     & 0.106 & 0.140 & 0.084 / 0.130 \\
\bottomrule
\end{tabular}
\end{table}

Table~\ref{tab:levels} breaks the results down by obstruction level. Two trends
emerge. First, distance MAE stays roughly constant ($\approx0.9$\,m) across all
obstruction levels, confirming that ToF-based ranging is largely immune to wall
loss once the path gain is accounted for. Second, angle MAE remains low
($7.6^{\circ}$--$9.8^{\circ}$) up to $20$\,dB wall loss, then climbs toward
$\approx90^{\circ}$---the expected MAE of a uniform random angle---for
obstruction $\ge40$\,dB. The confidence head tracks this transition: its MAE
and the mean predicted confidence fall monotonically as the wall loss grows,
meaning the model learns to express \emph{uncertainty} precisely where the angle
becomes uninformative.

\begin{table}[t]
\centering
\caption{Test performance per obstruction level.}
\label{tab:levels}
\begin{tabular}{ccccc}
\toprule
Wall loss (dB) & Confidence label & Dist MAE (m) & Angle MAE ($^{\circ}$) & Conf MAE \\
\midrule
0   & 1.0 & 0.607 & 8.20  & 0.032 \\
10  & 0.9 & 0.515 & 7.65  & 0.045 \\
20  & 0.8 & 0.708 & 9.83  & 0.051 \\
30  & 0.7 & 0.900 & 24.82 & 0.095 \\
40  & 0.6 & 0.907 & 67.35 & 0.194 \\
50  & 0.5 & 0.890 & 88.61 & 0.151 \\
60  & 0.4 & 0.886 & 90.65 & 0.061 \\
70  & 0.3 & 0.888 & 90.85 & 0.045 \\
80  & 0.2 & 0.883 & 90.89 & 0.145 \\
90  & 0.1 & 0.887 & 90.86 & 0.245 \\
\bottomrule
\end{tabular}
\end{table}

\subsection{Distributed Localization}
The full pipeline is evaluated on a $100\times100$\,m indoor scene containing
one BS and $15$ UEs partitioned by six $0.3$\,m-thick walls. Starting from the
BS, the confidence-greedy algorithm localizes nodes one at a time, each time
selecting the highest-confidence reachable link. Because NLOS links receive low
confidence, the localization tree preferentially grows through door-free LOS
links, and the graph refinement then corrects residual drift using all
high-confidence ($\bar{c}\ge0.8$) links jointly. The mean position error over
all UEs, reported before and after optimization, quantifies both the accuracy
of the sequential stage and the benefit of global refinement.

\section{Conclusion}
\label{sec:conclusion}

We introduced a unified framework for indoor localization that combines a
single-tower Transformer jointly estimating distance, angle, and confidence
from Wi-Fi CSI with a confidence-driven distributed multi-hop localization
algorithm. The confidence-weighted multi-task loss explicitly encodes the
physical reality that AoA information vanishes under NLOS obstruction, enabling
the model to learn \emph{when} to trust its own angle estimates. The resulting
confidence scores serve a dual purpose: they prioritize reliable links during
sequential localization and weight a subsequent graph optimization that removes
accumulated drift. Experiments on a multi-level wall-obstruction dataset
demonstrate sub-meter distance accuracy and single-digit degree angle accuracy
in LOS, with the confidence head cleanly discriminating LOS from NLOS links.
Future work will extend the approach to multi-AP scenarios and validate it on
real-world CSI traces.

\section*{Acknowledgment}
The authors would like to thank the anonymous reviewers for their constructive
comments.

\begin{thebibliography}{00}
\bibitem{halperin2011tool}
D.~Halperin, W.~Hu, A.~Sheth, and D.~Wetherall,
``Tool release: Gathering 802.11n traces with channel state information,''
\emph{ACM SIGCOMM Comput.\ Commun.\ Rev.}, vol.~41, no.~1, pp.~53--53, 2011.

\bibitem{kotaru2015spotfi}
M.~Kotaru, K.~Joshi, D.~Bharadia, and S.~Katti,
``SpotFi: Decimeter level localization using WiFi,''
in \emph{Proc.\ ACM SIGCOMM}, 2015, pp.~269--282.

\bibitem{xiong2013arraytrack}
J.~Xiong and K.~Jamieson,
``ArrayTrack: A fine-grained indoor location system,''
in \emph{Proc.\ USENIX NSDI}, 2013, pp.~71--84.

\bibitem{wang2015deepfi}
X.~Wang, L.~Gao, S.~Mao, and S.~Pandey,
``DeepFi: Deep learning for indoor fingerprinting using channel state
information,''
in \emph{Proc.\ IEEE WCNC}, 2015, pp.~1666--1671.

\bibitem{wang2016phasefi}
X.~Wang, L.~Gao, and S.~Mao,
``CSI phase fingerprinting for indoor localization with a deep learning
approach,''
\emph{IEEE Internet Things J.}, vol.~3, no.~6, pp.~1113--1123, 2016.

\bibitem{chen2017confi}
H.~Chen, Y.~Zhang, W.~Li, X.~Tao, and P.~Zhang,
``ConFi: Convolutional neural networks based indoor Wi-Fi localization using
channel state information,''
\emph{IEEE Access}, vol.~5, pp.~18066--18074, 2017.

\bibitem{vaswani2017attention}
A.~Vaswani et al.,
``Attention is all you need,''
in \emph{Proc.\ NeurIPS}, 2017, pp.~5998--6008.

\end{thebibliography}

\balance

\end{document}
